% Gerado por src/python/thesis/build_tab_tuning_ablacao.py -- nao editar a mao.
% Regenerar com: make thesis-tables
% Fonte: results/tuning_ivfspea2v2/head_to_head_c26_a43_v1_summary.json
\begin{table}[!htb]
    \centering
    \caption{Calibração de parâmetros: a configuração promovida frente à melhor configuração da fase inicial e à configuração inicial não ajustada.}
    \label{tab:tuning_c26}
    \resizebox{\textwidth}{!}{%
    \begin{tabular}{|c|l|l|c|}
        \hline
        \textbf{Métrica} & \textbf{Configuração} & \textbf{Comparador} & \textbf{V/E/D} \\
        \hline
        IGD & Configuração promovida & Melhor configuração da fase A & 0/12/0 \\
         & Configuração promovida & Configuração inicial não ajustada & 4/8/0 \\
         & Melhor configuração da fase A & Configuração inicial não ajustada & 4/8/0 \\
        \hline
        HV & Configuração promovida & Melhor configuração da fase A & 1/11/0 \\
         & Configuração promovida & Configuração inicial não ajustada & 4/7/1 \\
         & Melhor configuração da fase A & Configuração inicial não ajustada & 3/8/1 \\
        \hline
    \end{tabular}
    }
    \par\smallskip
    \parbox{0.95\textwidth}{\footnotesize
    Comparações par a par sobre as 12 configurações problema--objetivo do subconjunto de calibração, 30 execuções cada, orçamento de 50.000 avaliações. Teste de Mann--Whitney com correção de Holm, $\alpha = 0{,}05$. A configuração promovida não é distinguível da melhor configuração da fase A em nenhuma instância, e ambas superam a configuração inicial não ajustada em quatro delas: a calibração escolhe entre configurações equivalentes sem introduzir perda. Esta é evidência de justificativa da implementação, não da comparação principal.
    }
\end{table}
